\documentclass[12pt,a4paper]{article}
\usepackage[UTF8]{ctex}
\usepackage{amsmath,amssymb,amsthm}
\usepackage{geometry}

\geometry{margin=2.5cm}

\title{从任务空间wrench到关节力矩的转换}
\author{第8章补充说明}
\date{}

\begin{document}

\maketitle

\section{问题提出}

\textbf{问题}：要产生末端执行器加速度$\dot{V}$，需要施加wrench $F$，然后这个$F$再通过$J^T$去算到$\tau$吗？

\textbf{简短回答}：\textbf{是的，完全正确！}

\section{完整流程}

\subsection{步骤1：任务空间动力学}

\textbf{任务空间动力学方程}：
\begin{equation}
F = \Lambda(\theta)\dot{V} + \eta(\theta, V) \tag{8.89}
\end{equation}

\textbf{含义}：
\begin{itemize}
\item 要产生末端执行器加速度$\dot{V}$，需要施加wrench $F$
\item $F$包括两部分：
\begin{enumerate}
\item 惯性力：$\Lambda(\theta)\dot{V}$（产生加速度所需的力）
\item 科里奥利、向心力和重力：$\eta(\theta, V)$（由于速度和位置产生的力）
\end{enumerate}
\end{itemize}

\subsection{步骤2：从wrench到关节力矩}

\textbf{虚功原理}：
\begin{equation}
\tau = J^T(\theta)F \tag{8.51}
\end{equation}

\textbf{含义}：
\begin{itemize}
\item 任务空间wrench $F$在关节空间中的等效力矩是$\tau = J^T(\theta)F$
\item 这表示"产生相同的功率"
\end{itemize}

\subsection{完整流程总结}

\textbf{流程}：
\begin{enumerate}
\item \textbf{计算需要的wrench}：
\begin{equation}
F = \Lambda(\theta)\dot{V}_d + \eta(\theta, V)
\end{equation}
其中$\dot{V}_d$是期望的末端执行器加速度。

\item \textbf{转换为关节力矩}：
\begin{equation}
\tau = J^T(\theta)F
\end{equation}
\end{enumerate}

\textbf{完整公式}：
\begin{equation}
\tau = J^T(\theta)\left[\Lambda(\theta)\dot{V}_d + \eta(\theta, V)\right] \tag{*}
\end{equation}

\section{虚功原理的推导}

\subsection{功率守恒}

\textbf{虚功原理}告诉我们，关节空间和任务空间之间的功率是守恒的：
\begin{equation}
P_{\text{joint}} = P_{\text{task}}
\end{equation}

\textbf{关节空间功率}：
\begin{equation}
P_{\text{joint}} = \tau^T\dot{\theta}
\end{equation}

\textbf{任务空间功率}：
\begin{equation}
P_{\text{task}} = F^T V
\end{equation}

其中$V = J(\theta)\dot{\theta}$是末端执行器的twist。

\subsection{推导过程}

\textbf{步骤1：功率相等}
\begin{equation}
\tau^T\dot{\theta} = F^T V
\end{equation}

\textbf{步骤2：代入速度关系}
\begin{equation}
V = J(\theta)\dot{\theta}
\end{equation}

因此：
\begin{equation}
\tau^T\dot{\theta} = F^T J(\theta)\dot{\theta}
\end{equation}

\textbf{步骤3：提取$\dot{\theta}$}
\begin{equation}
(\tau^T - F^T J(\theta))\dot{\theta} = 0
\end{equation}

\textbf{步骤4：对所有$\dot{\theta}$成立}
由于这个关系对所有$\dot{\theta}$都成立，因此：
\begin{equation}
\tau^T = F^T J(\theta)
\end{equation}

\textbf{步骤5：转置}
\begin{equation}
\tau = J^T(\theta)F \tag{8.51}
\end{equation}

\section{实际应用：任务空间控制}

\subsection{控制目标}

\textbf{目标}：控制末端执行器跟踪期望轨迹

\textbf{期望量}：
\begin{itemize}
\item 期望位置：$X_d$
\item 期望速度：$V_d$
\item 期望加速度：$\dot{V}_d$
\end{itemize}

\subsection{控制律设计}

\textbf{步骤1：计算位置和速度误差}

\textbf{位置误差}（使用指数坐标）：
\begin{equation}
X_e = \log(X^{-1} X_d)
\end{equation}

\textbf{速度误差}：
\begin{equation}
V_e = [\text{Ad}_{X^{-1}X_d}]V_d - V
\end{equation}

\textbf{步骤2：计算期望加速度}

\textbf{PID控制}：
\begin{equation}
\dot{V}_d = \dot{V}_{d,\text{ref}} + K_p X_e + K_i \int X_e dt + K_d V_e
\end{equation}

其中$\dot{V}_{d,\text{ref}}$是参考加速度（前馈项）。

\textbf{步骤3：计算需要的wrench}

\textbf{使用任务空间动力学}：
\begin{equation}
F = \Lambda(\theta)\dot{V}_d + \eta(\theta, V) \tag{8.89}
\end{equation}

\textbf{或者添加反馈项}：
\begin{equation}
F = \Lambda(\theta)\dot{V}_d + \eta(\theta, V) + K_p X_e + K_d V_e
\end{equation}

\textbf{步骤4：转换为关节力矩}

\textbf{使用虚功原理}：
\begin{equation}
\tau = J^T(\theta)F \tag{8.51}
\end{equation}

\textbf{完整控制律}：
\begin{equation}
\tau = J^T(\theta)\left[\Lambda(\theta)\dot{V}_d + \eta(\theta, V) + K_p X_e + K_d V_e\right] \tag{11.47}
\end{equation}

\section{数学验证}

\subsection{从任务空间到关节空间}

\textbf{任务空间动力学}：
\begin{equation}
F = \Lambda(\theta)\dot{V} + \eta(\theta, V) \tag{8.89}
\end{equation}

\textbf{转换为关节空间}：
\begin{equation}
\tau = J^T(\theta)F = J^T(\theta)\left[\Lambda(\theta)\dot{V} + \eta(\theta, V)\right]
\end{equation}

\textbf{展开}：
\begin{align}
\tau &= J^T(\theta)\Lambda(\theta)\dot{V} + J^T(\theta)\eta(\theta, V) \\
&= J^T(\theta)J^{-T}(\theta)M(\theta)J^{-1}(\theta)\dot{V} + J^T(\theta)\eta(\theta, V) \\
&= M(\theta)J^{-1}(\theta)\dot{V} + J^T(\theta)\eta(\theta, V)
\end{align}

\textbf{代入$\dot{V} = \dot{J}(\theta)\dot{\theta} + J(\theta)\ddot{\theta}$}：
\begin{align}
\tau &= M(\theta)J^{-1}(\theta)(\dot{J}(\theta)\dot{\theta} + J(\theta)\ddot{\theta}) + J^T(\theta)\eta(\theta, V) \\
&= M(\theta)J^{-1}(\theta)\dot{J}(\theta)\dot{\theta} + M(\theta)\ddot{\theta} + J^T(\theta)\eta(\theta, V)
\end{align}

\textbf{代入$\eta(\theta, V)$的定义}：
\begin{align}
\eta(\theta, V) &= J^{-T}(\theta)h(\theta, J^{-1}(\theta)V) - \Lambda(\theta)\dot{J}(\theta)J^{-1}(\theta)V \\
&= J^{-T}(\theta)h(\theta, \dot{\theta}) - J^{-T}(\theta)M(\theta)J^{-1}(\theta)\dot{J}(\theta)J^{-1}(\theta)V
\end{align}

\textbf{代入$V = J(\theta)\dot{\theta}$}：
\begin{align}
\eta(\theta, V) &= J^{-T}(\theta)h(\theta, \dot{\theta}) - J^{-T}(\theta)M(\theta)J^{-1}(\theta)\dot{J}(\theta)\dot{\theta} \\
&= J^{-T}(\theta)h(\theta, \dot{\theta}) - J^{-T}(\theta)M(\theta)J^{-1}(\theta)\dot{J}(\theta)\dot{\theta}
\end{align}

\textbf{因此}：
\begin{align}
J^T(\theta)\eta(\theta, V) &= h(\theta, \dot{\theta}) - M(\theta)J^{-1}(\theta)\dot{J}(\theta)\dot{\theta}
\end{align}

\textbf{代入$\tau$的表达式}：
\begin{align}
\tau &= M(\theta)J^{-1}(\theta)\dot{J}(\theta)\dot{\theta} + M(\theta)\ddot{\theta} + h(\theta, \dot{\theta}) - M(\theta)J^{-1}(\theta)\dot{J}(\theta)\dot{\theta} \\
&= M(\theta)\ddot{\theta} + h(\theta, \dot{\theta})
\end{align}

\textbf{这正是关节空间动力学方程(8.82)！}

\section{物理意义}

\subsection{为什么需要这个转换？}

\textbf{任务空间控制的优势}：
\begin{itemize}
\item 直接控制末端执行器的运动
\item 更直观（我们关心的是末端执行器的位置，而不是关节角度）
\item 便于处理任务约束
\end{itemize}

\textbf{但执行器在关节空间}：
\begin{itemize}
\item 电机产生的是关节力矩$\tau$
\item 不能直接产生任务空间wrench $F$
\item 需要通过$J^T(\theta)$转换
\end{itemize}

\subsection{转换的物理意义}

\textbf{虚功原理的物理意义}：
\begin{itemize}
\item 关节空间和任务空间的功率是守恒的
\item 如果任务空间wrench $F$产生功率$F^T V$
\item 那么关节空间力矩$\tau$产生相同的功率$\tau^T\dot{\theta}$
\item 因此$\tau = J^T(\theta)F$
\end{itemize}

\textbf{具体例子}：

假设末端执行器需要产生力$F = [0, 0, 10, 0, 0, 0]^T$（沿$z$轴10 N的力）。

\textbf{步骤1}：计算需要的wrench
\begin{equation}
F = [0, 0, 10, 0, 0, 0]^T
\end{equation}

\textbf{步骤2}：转换为关节力矩
\begin{equation}
\tau = J^T(\theta)F
\end{equation}

\textbf{结果}：
\begin{itemize}
\item 每个关节的力矩$\tau_i$取决于雅可比矩阵$J(\theta)$
\item 如果关节$i$对$z$方向运动有贡献，那么$\tau_i \neq 0$
\item 如果关节$i$对$z$方向运动没有贡献，那么$\tau_i = 0$
\end{itemize}

\section{完整控制流程示例}

\subsection{任务空间轨迹跟踪}

\textbf{输入}：
\begin{itemize}
\item 期望轨迹：$X_d(t), V_d(t), \dot{V}_d(t)$
\item 当前状态：$X(t), V(t)$
\item 关节角度：$\theta(t)$
\end{itemize}

\textbf{步骤1：计算误差}
\begin{align}
X_e &= \log(X^{-1} X_d) \\
V_e &= [\text{Ad}_{X^{-1}X_d}]V_d - V
\end{align}

\textbf{步骤2：计算期望加速度}
\begin{equation}
\dot{V}_d = \dot{V}_{d,\text{ref}} + K_p X_e + K_d V_e
\end{equation}

\textbf{步骤3：计算需要的wrench}
\begin{equation}
F = \Lambda(\theta)\dot{V}_d + \eta(\theta, V) + K_p X_e + K_d V_e
\end{equation}

\textbf{步骤4：转换为关节力矩}
\begin{equation}
\tau = J^T(\theta)F
\end{equation}

\textbf{步骤5：执行控制}
\begin{itemize}
\item 将$\tau$发送给执行器
\item 执行器产生关节力矩
\item 机器人运动
\item 重复步骤1-5
\end{itemize}

\section{总结}

\subsection{关键要点}

\begin{enumerate}
\item \textbf{任务空间动力学}：
\begin{equation}
F = \Lambda(\theta)\dot{V} + \eta(\theta, V)
\end{equation}
告诉我们，要产生加速度$\dot{V}$，需要wrench $F$。

\item \textbf{虚功原理}：
\begin{equation}
\tau = J^T(\theta)F
\end{equation}
告诉我们，任务空间wrench $F$在关节空间中的等效力矩是$\tau$。

\item \textbf{完整流程}：
\begin{enumerate}
\item 计算需要的wrench：$F = \Lambda(\theta)\dot{V}_d + \eta(\theta, V)$
\item 转换为关节力矩：$\tau = J^T(\theta)F$
\end{enumerate}

\item \textbf{物理意义}：
\begin{itemize}
\item 任务空间控制更直观
\item 但执行器在关节空间
\item 需要通过$J^T(\theta)$转换
\item 转换基于功率守恒（虚功原理）
\end{itemize}
\end{enumerate}

\subsection{公式总结}

\textbf{任务空间动力学}：
\begin{equation}
F = \Lambda(\theta)\dot{V} + \eta(\theta, V) \tag{8.89}
\end{equation}

\textbf{虚功原理}：
\begin{equation}
\tau = J^T(\theta)F \tag{8.51}
\end{equation}

\textbf{完整控制律}：
\begin{equation}
\tau = J^T(\theta)\left[\Lambda(\theta)\dot{V}_d + \eta(\theta, V) + K_p X_e + K_d V_e\right] \tag{11.47}
\end{equation}

\end{document}

